% ============================================================
% BASELINE — Who&When adaptation for the MAST benchmark
%
% Two blocks:
%   (1) main-text snippet: 1-2 sentences for the body of the paper
%   (2) appendix block:    full methodological detail
%
% Source paper:
%   Yin et al., "Which Agent Causes Task Failures and When?
%   On Automated Failure Attribution of LLM Multi-Agent Systems"
%   (ICML 2025 Spotlight; arXiv:2505.00212)
%
% Adapted code lives in:
%   baselines/who&when/run_who_and_when_vllm.py
% ============================================================


% ---------- (1) MAIN-TEXT SNIPPET (1-2 sentences) -----------------

\paragraph{Who\&When baselines.}
We compare against the all-at-once (\textsc{W1}) and step-by-step
(\textsc{W2}) localization strategies of
\citet{yin2025whoandwhen}, adapted from single-agent attribution to
MAST's trace-level, 13-category multi-label setting. We omit the
binary-search variant (\textsc{W3}): once the single-error
assumption is dropped to recover MAST's multi-label outputs, the
procedure must run one bisection per category with both halves
allowed positive, which destroys the $\mathcal{O}(\log N)$
efficiency guarantee that motivated \textsc{W3} in the original
design.


% ---------- (2) APPENDIX BLOCK (full detail) ----------------------

\section{Who\&When Adaptation Details}
\label{app:whoandwhen_mast}

\citet{yin2025whoandwhen} introduce three localization strategies
for failure attribution in multi-agent systems: \textsc{W1}
(all-at-once), \textsc{W2} (step-by-step), and \textsc{W3} (binary
search). Each answers the same question --- \emph{which agent caused
the failure, and at which step?} --- under the assumption that
exactly one agent and one step are responsible per trace, returning
a single \texttt{(agent, step)} prediction. We adapt these
strategies to MAST, which is a \emph{trace-level} multi-label
benchmark over 13 failure modes
(1.1--1.5, 2.1--2.4, 2.6, 3.1--3.3): there is no per-step ground
truth, only a 13-bit indicator vector per trace
(\texttt{mast\_annotation}). The adapted runner is
\texttt{baselines/who\&when/run\_who\_and\_when\_vllm.py}; all
variants emit the existing yes/no output schema so the standard
scorer (\texttt{eval/calculate\_scores\_yesno.py}) is reused
without modification.

\paragraph{W1 --- All-at-Once.}
A single prompt presents the full trace together with the
definitions and examples of all 13 failure modes loaded from
\texttt{taxonomy\_definitions\_examples/}. The prompt mirrors the
original \textsc{W1} framing of \citet{yin2025whoandwhen}: ``You
are an AI assistant tasked with analyzing a multiagent system trace
when solving a real-world problem. The problem is:
\{task\_description\}. \ldots Identify which failure modes are
present, and explain the reason for each.'' The task description
is extracted from the first step (e.g.\ the \texttt{[mathproxyagent]
Problem: \ldots} preamble) and is included verbatim, matching
\citet{yin2025whoandwhen}'s use of \texttt{\{problem\}}; the
\texttt{ground\_truth} field used in the original is dropped
because MAST evaluation is reference-free. The output is a flat
13-bit yes/no vector (one bit per failure mode); MAST's existing
single-call yes/no baseline is structurally equivalent, so we
include this variant under the \textsc{W1} name primarily so the
three Who\&When variants share a common interface. One LLM call
per trace.

\paragraph{W2 --- Step-by-Step.}
For each step \texttt{step\_i} in the conversation, the model is
shown the cumulative history up to and including \texttt{step\_i}
and asked which of the 13 failure modes show direct evidence
\emph{at the current step}. The prompt body, the task description,
and the original's
``\emph{Note: avoid being overly critical \ldots focus on errors
that clearly derail the process}'' calibration sentence are
retained verbatim from \citet{yin2025whoandwhen}; we additionally
keep a step-locality clarification (``mark a mode only if the
current step itself provides evidence --- not a downstream
consequence of an earlier step'') that addresses the
trace-spanning blame ambiguity introduced by \textsc{W2}'s
sequential framing. The original \textsc{W2}'s early-exit on the
first ``yes'' is removed; we sweep the full trace and
OR-aggregate the per-step 13-bit vectors into a single
trace-level prediction. This re-architecting is necessary because
the trace-level ground truth requires the union of evidence
across all steps, not just the first occurrence. To keep the
prompt within the model's context budget we cap the cumulative
history block at $80{,}000$ characters with sliding-window
truncation: the oldest steps are dropped first and replaced by an
explicit \texttt{[\ldots\ N earlier step(s) omitted \ldots]}
marker, so the most recent context (and the current step) are
always preserved. The threshold rarely binds on AG2 traces.
$N$ LLM calls per trace, where $N$ is the number of steps in
\texttt{steps[]} (median $\approx 5$--$15$ on AG2 traces, with
long tails).

\paragraph{Why we omit W3.}
The original \textsc{W3} is binary search: at each level the model
chooses between the upper and lower half (exactly one), and
recursion continues only into the chosen side until a single
``decisive'' step is isolated. This procedure produces \emph{one}
prediction in $\mathcal{O}(\log N)$ calls --- the algorithmic
motivation for \textsc{W3} in the single-error setting.
MAST is multi-label: any of the 13 failure modes may be present
or absent independently, and a present mode may appear in any
subset of segments (including both halves of any interval). A
faithful multi-label adaptation must therefore (i) run a separate
bisection per failure mode and (ii) reformulate the per-level
question as two independent booleans
(\textit{lower\_half\_present}, \textit{upper\_half\_present})
rather than a forced choice. The first change multiplies the call
count by $|\mathcal{C}|=13$; the second eliminates the
binary-search efficiency guarantee, because in the worst case (a
label present everywhere) the recursion enumerates every leaf for
$\mathcal{O}(N)$ calls per label and $\mathcal{O}(|\mathcal{C}|
\cdot N)$ calls overall --- strictly more expensive than
\textsc{W2}'s $\mathcal{O}(N)$ sweep, with no remaining
sub-linear localization advantage.
For MAST the gap is even more pronounced than for span-level
benchmarks: because the label is trace-level, the only signal
\textsc{W3} can hope to extract is ``does the mode occur
\emph{anywhere}'', which a single \textsc{W1} call already
provides at $1/13\!\times\!\log_2 N$ of the cost. We therefore
include \textsc{W1} and \textsc{W2} as faithful adaptations of
the Who\&When framework, and exclude \textsc{W3} on the grounds
that its design hypothesis (sub-linear localization of a unique
error) is incompatible with MAST's multi-label, trace-level
ground truth.

\paragraph{Implementation notes.}
Both variants share the same response parser, which strips
reasoning tags (\texttt{<think>...</think>}, including the
truncated-open case where the closing tag is missing because the
model hit its token budget mid-thought) and gpt-oss Harmony
channel markers
(\texttt{assistantanalysis}/\texttt{commentary}/\texttt{final},
keeping only the \texttt{assistantfinal} payload) before
extracting the per-mode yes/no vector with regex matching on
\texttt{<mode\_id>: yes|no}. For reasoning models
(QwenLong-L1-32B, gpt-oss, DeepSeek-R1, QwQ) the runner
auto-bumps \texttt{max\_tokens} from 2000 to 16000, because the
hidden reasoning chain consumes most of the default budget
before the answer block and would otherwise silently default
the entire 13-bit vector to zero (a false-negative rather than
a parse failure). The \textsc{W2} runner additionally caches
per-step predictions in the output \texttt{meta.per\_step}
field so post-hoc analyses can reconstruct the OR-aggregation
across steps.

\paragraph{Faithfulness audit: divergences from the original prompts.}
In prompting-based LLM-as-judge the \emph{prompt is the method}, so we
catalogue every textual change between the original prompts of
\citet{yin2025whoandwhen} (\texttt{Automated\_FA/Lib/utils.py}) and our
MAST adaptation. Each change is either \textbf{forced} (an unavoidable
consequence of the single$\to$multi remap and MAST's trace-level,
13-bit, reference-free output schema) or a \textbf{deviation} (a choice
that goes beyond that remap). After the corrections described in this
section, \textsc{W1} and \textsc{W2} retain only forced changes plus a
single MAST-side adapter (the \textsc{W2} cumulative-history truncation,
which the original Who\&When does not need because its chat histories
are short).

\begin{itemize}\setlength{\itemsep}{0.4em}

\item \noindent\rule{\linewidth}{0.4pt}\\
      \textbf{Aspect:}\enspace \textsc{W1}/\textsc{W2}: 13-mode failure-mode taxonomy + definitions + examples\\
      \textbf{Original W\&W:}\enspace absent (free-form ``identify the agent and the step'')\\
      \textbf{Our adaptation:}\enspace full taxonomy with definitions and examples injected at prompt head\\
      \textbf{Justified by single$\to$multi?:}\enspace \textbf{forced} --- the 13 modes are the multi-label output space\\
      \rule{\linewidth}{0.4pt}

\item \textbf{Aspect:}\enspace \textsc{W1}/\textsc{W2}: ``who'' (which agent) $\to$ ``what'' (which failure mode); ``when'' axis dropped\\
      \textbf{Original W\&W:}\enspace predicts a single \texttt{(agent, step, reason)} tuple\\
      \textbf{Our adaptation:}\enspace predicts a 13-bit yes/no vector at the trace level; no step output\\
      \textbf{Justified by single$\to$multi?:}\enspace \textbf{forced} --- MAST's ground truth is the trace-level \texttt{mast\_annotation} 13-bit vector; there is no per-step gold label\\
      \rule{\linewidth}{0.4pt}

\item \textbf{Aspect:}\enspace \textsc{W1}/\textsc{W2}: ground-truth answer (\texttt{ground\_truth})\\
      \textbf{Original W\&W:}\enspace included; the original authors annotate that it ``should be removed if ground truth shouldn't be used''\\
      \textbf{Our adaptation:}\enspace dropped\\
      \textbf{Justified by single$\to$multi?:}\enspace \textbf{forced} --- MAST evaluation is reference-free; including it would be data leakage\\
      \rule{\linewidth}{0.4pt}

\item \textbf{Aspect:}\enspace \textsc{W1}/\textsc{W2}: output format\\
      \textbf{Original W\&W:}\enspace free text (\texttt{Agent Name: \ldots Step Number: \ldots Reason: \ldots} or \texttt{1. Yes/No 2. Reason: \ldots})\\
      \textbf{Our adaptation:}\enspace \texttt{@@ \ldots @@}-delimited block with one \texttt{<mode\_id>: yes/no} line per failure mode, parsed by regex\\
      \textbf{Justified by single$\to$multi?:}\enspace \textbf{forced} --- 13 binary labels need a structured output; the existing MAST scorer ingests this format\\
      \rule{\linewidth}{0.4pt}

\item \textbf{Aspect:}\enspace \textsc{W2}: early exit on first ``yes''\\
      \textbf{Original W\&W:}\enspace defining behaviour --- the loop terminates as soon as one error is localized\\
      \textbf{Our adaptation:}\enspace removed; the sweep covers every step and the per-step 13-bit vectors are OR-aggregated into a trace-level prediction\\
      \textbf{Justified by single$\to$multi?:}\enspace \textbf{forced} --- a single-error early exit cannot recover the union of failure modes occurring across the trace\\
      \rule{\linewidth}{0.4pt}

\item \textbf{Aspect:}\enspace \textsc{W2}: cumulative-history truncation\\
      \textbf{Original W\&W:}\enspace no truncation --- relies on a reactive Azure OpenAI failure if the prompt is too long\\
      \textbf{Our adaptation:}\enspace sliding-window cap at 80{,}000 cumulative-history characters; oldest steps dropped first and replaced by an explicit \texttt{[\ldots\ N earlier step(s) omitted \ldots]} marker\\
      \textbf{Justified by single$\to$multi?:}\enspace \textbf{deviation} (MAST-side adapter) --- MAST traces have a long tail that the original's reactive bail-out handles unreliably under vLLM; the cap is honest about what it drops and rarely binds\\
      \rule{\linewidth}{0.4pt}

\end{itemize}

The earlier draft of this baseline contained three additional deviations
(a dropped task description in both \textsc{W1} and \textsc{W2}, an
aggressive ``for EACH of the 13 modes, INDEPENDENT decision'' framing
in \textsc{W1}, and a missing
``\emph{avoid being overly critical}'' calibration sentence in
\textsc{W2}). All three have been corrected in the runner
(\texttt{baselines/who\&when/run\_who\_and\_when\_vllm.py}); the
previous version is preserved at
\texttt{old\_run\_who\_and\_when\_vllm.py} for diffing. The remaining
changes are forced by the multi-label, trace-level adaptation, with
the cumulative-history truncation as the only non-forced MAST-side
adapter.
